\section*{第 9 讲\quad 希尔伯特定理 90 与上链复形}

一如既往，$G=\Z/n\Z$，$L/K$ 是局部域的 Galois 扩张，$\Gal(L/K)=G$，生成元为 $\sigma\in G$。上一讲我们证明了：

\textbf{定理 9.1.} $\chi(L^\times)=n$，其中 $\chi$ 表示 Herbrand 商。

注意，我们的方法其实很容易推广到非阿基米德的情形。本讲我们将证明：

\textbf{定理 9.2（希尔伯特定理 90）.} $\HH^1(G,L^\times)=0$。\footnote{原文此处印作 $\HH^1(G,K^\times)=0$；由下文 $\HH^0(G,L^\times)=K^\times/N(L^\times)$、以及推论 9.3 给出的显式陈述可知应为 $\HH^1(G,L^\times)=0$，此处径改。}

两者合起来便推出 $\HH^0(G,L^\times)=K^\times/N(L^\times)$ 的基数为 $n$。希尔伯特定理 90 的另一个推论是：若 $L$ 与 $K$ 都是有限域，则对一切 $i$ 有 $\HH^i(G,L^\times)=0$，因为 $L$ 有限故 $\chi(L^\times)=1$（于是由周期性所有上同调都消没）。这与 [Wei74] 中的第一个结果类似。具体地，
\[
\HH^1(G,L^\times)=\ker(N:L^\times\to K^\times)\big/\{y/\sigma y:y\in L^\times\},
\]
其中每个 $y/\sigma y$ 的范数都是 $1$，因为 $y/\sigma y\cdot\sigma y/\sigma^2y\cdots\sigma^{n-1}y/\sigma^ny=1$。于是希尔伯特定理 90 蕴含：

\textbf{推论 9.3.} 若 $L/K$ 是循环扩张，$x\in L^\times$ 且 $N(x)=1$，则存在 $y\in L^\times$ 使得 $x=y/\sigma y$。

\textbf{例 9.4.} 取 $L:=\Q(i)$、$K:=\Q$。选 $x\in\Q(i)$ 使 $N(x)=1$。则对某些满足 $a^2+b^2=c^2$ 的 $a,b,c\in\Z$ 有 $x=a/c+(b/c)i$。于是希尔伯特定理 90 给出勾股数组的通常参数化：
\[
(r-s)^2+(2rs)^2=(r+s)^2 .
\]

对 $n=2$，证明很简单。此时 $N(x)=x\cdot\sigma x=1$。若 $x\neq-1$，令 $y:=x+1$，则
\[
x\cdot\sigma y=x(\sigma x+1)=N(x)+x=1+x=y,
\]
故 $x=y/\sigma y$，如所欲证。若 $x=-1$，令 $y:=\sqrt d$，其中 $L=K(\sqrt d)$，则同样有 $y/\sigma y=\sqrt d/(-\sqrt d)=-1=x$。注意，这就完成了对局部域的二次扩张 $L/K$ 证明 $\#(K^\times/NL^\times)=2$，从而完成了希尔伯特符号良好性质的证明！事实上，回忆对域 $L:=K(\sqrt a)$（$a\in K^\times$ 但非平方元），$(a,b)=1$ 当且仅当 $b\in N(L^\times)$。

现在我们转向希尔伯特定理 90 的一般情形。主要引理如下：

\textbf{引理 9.5.} 对每个 $x\in L$，令
\[
H_x:L\to L,\qquad y\mapsto x\cdot\sigma(y),
\]
它是 $K$-向量空间的线性映射。则 $H_x$ 的特征多项式为 $t^n-N(x)\in K[t]$，其中我们把特征多项式的定义规范化为首一的。

注意，这个特征多项式比 $y\mapsto xy$ 的特征多项式简单：只要 $T(x)\neq 0$，后者就会含有 $t^n-1$ 的非零倍，而迹非退化（任何可分扩张都如此）时就会出现这种情况。

\textit{证明（由 9.5 推出 9.2）.} 取 $x\in L$ 并设 $N(x)=1$。则 $H_x$ 的特征多项式为 $t^n-1$，故 $1$ 是它的根，从而是 $H_x$ 的特征值。于是 $\ker((H_x-1)\otimes_K\bar K)\neq 0$；而对域而言张量积与取核交换，故 $\ker(H_x-1)\neq 0$。于是存在 $y\in L^\times$ 使得 $H_x(y)=x\cdot\sigma(y)=y$，即 $x=y/\sigma y$，如所欲证。\hfill$\square$

\textit{证明（引理 9.5）.} 首先注意 $H_x^n$ 对应乘以 $N(x)$，因为对任意 $y\in L$，
\[
H_x^n(y)=x\cdot\sigma(x\cdot\sigma(x\cdots\sigma(y)))=x\cdot\sigma(x)\cdot\sigma^2(x)\cdots\sigma^{n-1}(x)\cdot\sigma^n(y)=N(x)y .
\]
于是 $H_x$ 的极小多项式整除 $t^n-N(x)$。现在回忆：线性算子 $T$ 的极小多项式总整除其特征多项式，而后者次数为 $n$，故只需证明两者相等。这等价于 $T$ 的若尔当分解中没有共享特征值的块，而这又等价于对一切 $\lambda\in K$ 有 $\dim_K(\ker(T-\lambda I))\le 1$。

下面给出一个并不完全成立的证明。设 $H_x(y_1)=\lambda y_1$、$H_x(y_2)=\lambda y_2$，且 $y_1,y_2\neq 0$（即两者都是「真正的」特征向量）。我们想证明 $y_2$ 是 $y_1$ 的倍数，即 $y_2/y_1\in K$，也就是被 $\Gal(L/K)$ 固定。事实上，由 $\sigma y_2=H_x(y_2)/x$（对 $y_1$ 同理）可得
\[
\frac{\sigma y_2}{\sigma y_1}=\frac{x\lambda y_2}{x\lambda y_1}=\frac{y_2}{y_1}.
\]
然而问题在于：这个证明发生在 $L$ 中，而不在 $K$ 中——而我们的特征值实际住在 $K$ 里！因此，我们需要在 $L\otimes_K\bar K\simeq\prod_{g\in G}\bar K$ 中工作，而它不一定是域。

我们可以在基变换之后再计算特征多项式。回忆
\[
L\otimes_K\bar K\xrightarrow{\sim}\prod_{i=0}^{n-1}\bar K,\qquad a\otimes b\mapsto\big((\sigma^ia)\cdot b\big)_{i=0}^{n-1}.
\]
这个同构非典范地延拓为 $\bar K$ 的一个自同构，除此之外一切都是典范的，其中群通过左乘作用于坐标集合。映射
\[
\sigma\otimes\mathrm{id}:L\otimes_K\bar K\to L\otimes_K\bar K
\]
对应置换坐标；此外还有映射
\[
\mu_x\otimes\mathrm{id}:L\otimes_K\bar K\to L\otimes_K\bar K,\qquad
(y_0,\dots,y_{n-1})\mapsto\big(xy_0,(\sigma x)y_1,\dots,(\sigma^{n-1}x)y_{n-1}\big),
\]
其中 $\mu_x$ 表示乘以 $x$。现在设 $\lambda\in\bar K$ 是 $H_x$ 的特征值，对应的特征向量为 $(y_0,\dots,y_{n-1})$。则
\[
H_x(y)=\big(xy_1,(\sigma x)y_2,(\sigma^2x)y_3,\dots,(\sigma^{n-1}x)y_0\big)
=(\lambda y_0,\lambda y_1,\lambda y_2,\dots,\lambda y_{n-1}),
\]
于是 $xy_1=\lambda y_0$，即 $y_1=(\lambda/x)y_0$；同理 $y_2=(\lambda/\sigma x)y_1=(\lambda^2/(x\cdot\sigma x))y_0$。一般地有
\[
y_i=\frac{\lambda^i}{x\cdots\sigma^{i-1}x}\,y_0,
\]
故所有坐标都由 $y_0$ 唯一确定，即
\[
(y_0,\dots,y_{n-1})=y_0\Big(1,\frac{\lambda}{x},\frac{\lambda^2}{x\cdot\sigma x},\dots,\frac{\lambda^{n-1}}{\prod_{i=0}^{n-2}\sigma^ix}\Big).
\]
因此每个特征空间确实都是一维的，如所欲证。注意，只有当
\[
\frac{\lambda^n}{x\cdots\sigma^{n-1}x}=\frac{\lambda^n}{N(x)}=1,
\]
即 $\lambda^n=N(x)$ 时，这才真正定义出特征向量，这与我们的预期一致（并且 $n$ 个 $n$ 次根都以重数一出现）。\hfill$\square$

现在回忆我们的目标是：对局部域的交换扩张 $L/K$，典范地（且是在很强的意义下）证明
\[
K^\times/NL^\times\simeq\Gal(L/K).
\]
我们已经证明 $K^\times/NL^\times$ 的阶正确，但将用上同调对非循环群一般地证明它。下面我们引入同调代数的语言，这将是我们的核心工具。

\textbf{定义 9.6.} 交换群的（上链）复形 $X$ 是一列
\[
\cdots\to X^{-1}\xrightarrow{d^{-1}}X^0\xrightarrow{d^0}X^1\xrightarrow{d^1}\cdots,
\]
使得微分满足对每个 $i$ 有 $d^{i+1}d^i=0$。

\textbf{记号 9.7.} 我们常把整个复形记作 $X^\bullet$，其中「$\bullet$」处于指标的位置。我们也常省略指标，例如把 $d^i$ 写成 $d$，并写 $d\cdot d=d^2=0$。注意，有些作者写作 $H_i:=H^{-i}$，对 $X^i$ 也类似，这样微分降低次数；而我们的约定是微分提升次数。

\textbf{定义 9.8.} 第 $i$ 个上同调群为 $H^i(X):=\ker(d^i)/\im(d^{i-1})$。

这些上同调群其实正是我们要的不变量；但 $X$ 是更「丰富」的对象，因此最好在我们过程的最后一步才过渡到上同调。下面引入复形映射的零伦这一重要想法。

\textbf{定义 9.9.} 若下图交换
\[
\begin{CD}
\cdots @>>> X_{-1} @>d_{-1}>> X_0 @>d_0>> X_1 @>>> \cdots\\
@. @V{f_{-1}}VV @V{f_0}VV @V{f_1}VV @.\\
\cdots @>>> Y_{-1} @>d_{-1}>> Y_0 @>d_0>> Y_1 @>>> \cdots
\end{CD}
\]
则称 $f$ 为\textbf{复形映射}。注意，$f$ 诱导上同调之间的映射，因为由交换性，$X_\bullet$ 中微分的核与像都被保持到 $Y_\bullet$ 中。若下图中的映射 $h$
\[
\begin{CD}
\cdots @>>> X_{-1} @>d_{-1}>> X_0 @>d_0>> X_1 @>>> \cdots\\
@. @V{f_{-1}}VV @V{f_0}VV @V{f_1}VV @.\\
\cdots @>>> Y_{-1} @>d_{-1}>> Y_0 @>d_0>> Y_1 @>>> \cdots
\end{CD}
\]
（图中另有斜向连接相邻两行的映射 $h_0,h_1$，此处从略）满足 $dh+hd=f$，更精确地说，对每个 $i$ 有
\[
% [原文疑误：原印 $d_ih_{i+1}+h_id_{i+1}=f_{i+1}$ 指标位移有误，与引理 9.10 的证明不自洽；已按标准约定改写，其中 $h_i:X_i\to Y_{i-1}$]
f_i=d_{i-1}h_i+h_{i+1}d_i
\]
则称 $h$ 为 $f$ 的\textbf{零伦}。

\textbf{引理 9.10.} 若 $f$ 是零伦的，则其诱导的上同调映射 $H^i(X)\xrightarrow{H^i(f)}H^i(Y)$ 对一切 $i$ 为零。

\textit{证明.} 取 $x\in X_i$ 使得 $d_ix=0$。则 $f(x)=d_{i-1}h_i(x)+h_{i+1}d_ix=d(h_i(x))$，故 $f(x)\in\im(d_{i-1})$，从而 $f(x)=0$ 于 $H^i(Y)$。\hfill$\square$

现在，我们在此处的指导原则是：对代数而言，同构是比相等好得多的概念，因为相等指的是没有结构的集合。因此，若 $f\simeq g$，即 $f$ 同伦于 $g$（意思是 $f-g$ 存在零伦），则任何真正的数学检验都无法再区分 $f$ 与 $g$。

我们希望为复形映射定义某种「余核」的概念。一个糟糕的想法是：对复形映射 $f:X\to Y$ 直接取 $\coker(f)$。更好的想法如下：

\textbf{定义 9.11.} \textbf{同伦余核}或\textbf{锥} $h\coker(f)=\operatorname{Cone}(f)$ 具有如下泛性质：复形映射 $h\coker(f)\to Z$ 等价于映射 $Y\to Z$ 连同 $X\to Z$ 的一个零伦的数据；注意这给出如下交换图
\[
\begin{CD}
X @>>> Z\\
@V{f}VV @.\\
Y @.
\end{CD}
\]
注意它与通常余核的泛性质之间的相似性。
